\documentclass{../../note}

\usepackage{xcolor}
\usepackage{longtable}
\usepackage{array}
\usepackage{listings}

\lstset{
  basicstyle=\ttfamily\small,
  keywordstyle=\color{blue},
  commentstyle=\color{green!50!black},
  stringstyle=\color{red!60!black},
  numbers=left,
  numberstyle=\tiny,
  numbersep=6pt,
  breaklines=true,
  frame=single
}

\title{HUFE 数据库第二讲}
\author{isomoes}

\begin{document}

\maketitle

\section{本次课定位}

本次课对应 5 次课安排中的第 5 次课，时长 3 小时，内容覆盖\textbf{数据库安全性、数据库完整性、关系数据理论、数据库设计以及数据库恢复与并发控制}。这是数据库部分的收尾讲，也是整个课程的综合串讲。

这一讲的特点是：\textbf{概念分散但考点集中}。不同于第一讲可以用"系统—模型—操作"串起来，本讲更像五个相对独立的模块，需要分别抓住各模块的核心考点，再在最后串联成整体认知。

\subsection{本次课目标}

\begin{itemize}
  \item 理解数据库安全性的基本概念，掌握自主存取控制中 GRANT 和 REVOKE 的基本用法。
  \item 掌握实体完整性、参照完整性和用户定义完整性的定义与约束规则。
  \item 理解函数依赖的定义，掌握 1NF、2NF、3NF 和 BCNF 的判断方法。
  \item 掌握数据库设计的六个基本步骤，能绘制简单 E-R 图并将其转换为关系模式。
  \item 理解事务的 ACID 属性，了解故障类型与恢复技术的基本思路，掌握并发控制的基本概念。
\end{itemize}

\subsection{考试导向提醒}

本讲在试卷中属于\textbf{分散高频区}，考法灵活，常见方式包括：

\begin{itemize}
  \item \textbf{概念判断题}：如 ACID 属性、事务定义、安全性控制手段、完整性约束类型。
  \item \textbf{范式判断题}：给出关系模式和函数依赖集，判断属于第几范式并说明原因。
  \item \textbf{模式分解题}：将低范式关系模式分解成满足更高范式的模式组合。
  \item \textbf{E-R 图与关系模式转换题}：根据描述设计 E-R 图，再转换成关系模式。
\end{itemize}

这一讲的学习策略是：\textbf{每个模块各抓一条主线，不要混在一起背}。

\subsection{3小时建议节奏}

\begin{center}
\begin{tabular}{|c|c|p{9cm}|}
\hline
阶段 & 时间 & 主要内容 \\
\hline
第1段 & 30分钟 & 数据库安全性：基本概念、自主存取控制、GRANT/REVOKE \\
\hline
第2段 & 30分钟 & 数据库完整性：实体完整性、参照完整性、用户定义完整性 \\
\hline
第3段 & 45分钟 & 关系数据理论：函数依赖、范式（1NF～BCNF）、模式分解 \\
\hline
第4段 & 30分钟 & 数据库设计：六阶段、E-R 模型、关系模式转换 \\
\hline
第5段 & 25分钟 & 恢复与并发控制：事务 ACID、故障类型、恢复技术、并发问题 \\
\hline
\end{tabular}
\end{center}

\section{数据库安全性}

\subsection{为什么要有安全性控制}

数据库中往往存储着大量敏感数据。安全性控制的目标是：\textbf{确保只有合法的用户以合法的方式访问数据}。

数据库安全性涉及多个层面，但考试主要关注 DBMS 层面的安全机制。

\subsection{自主存取控制}

最常见的数据库安全控制机制是\textbf{自主存取控制（DAC）}，其核心是"谁拥有数据，谁决定谁能访问"。

在 SQL 中通过两条语句实现：

\begin{itemize}
  \item \texttt{GRANT}：授予权限；
  \item \texttt{REVOKE}：回收权限。
\end{itemize}

\subsubsection{GRANT 授权}

\begin{lstlisting}[language=SQL]
GRANT SELECT, INSERT
ON Student
TO user1
WITH GRANT OPTION;
\end{lstlisting}

关键点：

\begin{itemize}
  \item \texttt{ON} 后跟被授权对象，如表名；
  \item \texttt{TO} 后跟被授权用户；
  \item \texttt{WITH GRANT OPTION} 表示允许 user1 再把此权限转授他人。
\end{itemize}

\subsubsection{REVOKE 回收}

\begin{lstlisting}[language=SQL]
REVOKE SELECT
ON Student
FROM user1
CASCADE;
\end{lstlisting}

\texttt{CASCADE} 表示同时回收 user1 转授给他人的相应权限。

\subsection{强制存取控制}

强制存取控制（MAC）基于安全标签，由系统强制执行，用户不能自行修改。常见于对安全要求极高的系统。考试中只需了解其概念，区分 DAC 和 MAC 即可。

\subsection{数据库安全性练习}

\begin{enumerate}
  \item 自主存取控制和强制存取控制的主要区别是什么？
  \item \texttt{GRANT SELECT ON Student TO user1;} 语句的作用是什么？
  \item \texttt{WITH GRANT OPTION} 表示什么含义？
  \item 若要回收 user1 对 Student 表的查询权限，应如何写 SQL 语句？
  \item 请判断：自主存取控制中，数据对象的访问权限由 DBMS 强制决定，用户不能干预。（\hspace{1cm}）
  \item 请判断：使用 \texttt{CASCADE} 选项回收权限时，仅回收目标用户自己的权限，不影响其转授给他人的权限。（\hspace{1cm}）
\end{enumerate}

\section{数据库完整性}

\subsection{完整性的意义}

数据库完整性是指数据库中数据的\textbf{正确性、有效性和相容性}。它通过约束规则来防止不合语义的数据进入数据库。

完整性约束分为三类，这三类必须分清：

\begin{enumerate}
  \item \textbf{实体完整性}
  \item \textbf{参照完整性}
  \item \textbf{用户定义完整性}
\end{enumerate}

\subsection{实体完整性}

\textcolor{red}{实体完整性规则}：关系的主键属性值不能为空（NULL）。

这是因为主键的作用是唯一标识关系中的每一个元组，如果主键为空，就无法进行唯一标识。

在 SQL 中，定义主键时会自动施加此约束：

\begin{lstlisting}[language=SQL]
CREATE TABLE Student (
    Sno CHAR(9) PRIMARY KEY,
    Sname VARCHAR(20) NOT NULL,
    Sdept VARCHAR(20)
);
\end{lstlisting}

\subsection{参照完整性}

\textcolor{red}{参照完整性规则}：如果关系 $R$ 的外键 $F$ 是关系 $S$ 的主键，则 $R$ 中每个元组在 $F$ 上的取值必须：

\begin{itemize}
  \item 等于 $S$ 中某个元组的主键值；
  \item 或为空值（如果允许外键取空值）。
\end{itemize}

例如选课表 $SC(Sno, Cno, Grade)$ 中，$Sno$ 是外键，引用 $Student$ 表的主键。若 $Student$ 表中不存在某个 $Sno$，则 $SC$ 中不能插入含该 $Sno$ 的记录。

在 SQL 中定义外键约束：

\begin{lstlisting}[language=SQL]
CREATE TABLE SC (
    Sno CHAR(9),
    Cno CHAR(4),
    Grade SMALLINT,
    PRIMARY KEY (Sno, Cno),
    FOREIGN KEY (Sno) REFERENCES Student(Sno),
    FOREIGN KEY (Cno) REFERENCES Course(Cno)
);
\end{lstlisting}

\subsection{用户定义完整性}

用户定义完整性是针对具体应用语义的约束，例如：

\begin{itemize}
  \item 成绩必须在 0 到 100 之间；
  \item 性别只能是"男"或"女"。
\end{itemize}

在 SQL 中常用 \texttt{CHECK} 子句实现：

\begin{lstlisting}[language=SQL]
CREATE TABLE SC (
    Sno CHAR(9),
    Grade SMALLINT CHECK (Grade >= 0 AND Grade <= 100)
);
\end{lstlisting}

\subsection{完整性约束命名}

可以使用 \texttt{CONSTRAINT} 给约束命名，便于后续修改或删除：

\begin{lstlisting}[language=SQL]
CREATE TABLE SC (
    Grade SMALLINT,
    CONSTRAINT GradeCheck CHECK (Grade >= 0 AND Grade <= 100)
);
\end{lstlisting}

\subsection{数据库完整性练习}

\begin{enumerate}
  \item 什么是实体完整性？为什么主键不能为空？
  \item 参照完整性规则约束的是哪一类属性？
  \item 外键的取值有哪些合法情况？
  \item 用户定义完整性和前两类完整性的主要区别是什么？
  \item \texttt{CHECK} 子句在 SQL 中的作用是什么？请举一个例子。
  \item 请判断：外键的取值只能是被参照关系中已有的主键值，不允许为 NULL。（\hspace{1cm}）
\end{enumerate}

\section{关系数据理论}

\subsection{为什么要有规范化理论}

设计关系数据库时，如果关系模式设计不合理，就会出现：

\begin{itemize}
  \item \textbf{数据冗余}：同样的数据被反复存储；
  \item \textbf{更新异常}：修改一处可能要修改多处；
  \item \textbf{插入异常}：无法单独插入某些数据；
  \item \textbf{删除异常}：删除某数据时意外丢失其他信息。
\end{itemize}

规范化理论的目标就是通过合理分解关系模式，消除这些问题。

\subsection{函数依赖}

\subsubsection{定义}

设关系模式 $R(U)$，$X, Y \subseteq U$。若对 $R$ 的任意两个元组 $t_1, t_2$，只要 $t_1[X] = t_2[X]$，就有 $t_1[Y] = t_2[Y]$，则称 $X$ \textcolor{red}{函数确定} $Y$，记作 $X \to Y$，即 $Y$ \textbf{函数依赖}于 $X$。

通俗理解：\textbf{$X$ 的值确定之后，$Y$ 的值也就唯一确定了}。

例如：
\begin{itemize}
  \item 学号 $\to$ 学生姓名（已知学号，姓名唯一确定）；
  \item 学号 $\to$ 所在院系。
\end{itemize}

\subsubsection{完全函数依赖与部分函数依赖}

\begin{itemize}
  \item \textbf{完全函数依赖}：$X \to Y$，且 $X$ 的任何真子集都不能函数确定 $Y$。记作 $X \xrightarrow{F} Y$。
  \item \textbf{部分函数依赖}：$X \to Y$，但存在 $X$ 的某个真子集 $X'$ 使得 $X' \to Y$。记作 $X \xrightarrow{P} Y$。
\end{itemize}

例如：选课关系 $SC(Sno, Cno, Grade)$，$(Sno, Cno) \to Grade$ 是完全函数依赖，因为仅用学号或仅用课程号都不能唯一确定成绩。

\subsubsection{传递函数依赖}

若 $X \to Y$，$Y \to Z$，且 $Y \not\to X$，则称 $Z$ 传递函数依赖于 $X$。

例如：如果学号确定系名，系名确定系主任，则系主任传递依赖于学号。

\subsection{第一范式（1NF）}

\textcolor{red}{1NF}：关系模式 $R$ 的每个属性都是不可再分的原子值。

这是最基本的要求，所有关系数据库中的关系都必须满足 1NF。

判断要点：如果某一属性可以继续拆分成多个子属性，那么违反 1NF。

\subsection{第二范式（2NF）}

\textcolor{red}{2NF}：在满足 1NF 的基础上，每个非主属性\textbf{完全函数依赖}于主键（候选键）。

换句话说，若存在非主属性对主键的\textbf{部分函数依赖}，则不满足 2NF。

例如：若关系模式为 $R(Sno, Cno, Sname, Grade)$，其中 $(Sno, Cno)$ 是主键，但 $Sno \to Sname$（学号已能确定姓名），$Sname$ 部分依赖于主键，则 $R$ 不满足 2NF。

解决方法：分解为：
\begin{itemize}
  \item $Student(Sno, Sname)$；
  \item $SC(Sno, Cno, Grade)$。
\end{itemize}

\subsection{第三范式（3NF）}

\textcolor{red}{3NF}：在满足 2NF 的基础上，每个非主属性\textbf{不传递依赖}于主键。

若存在非主属性对主键的\textbf{传递函数依赖}，则不满足 3NF。

例如：若 $Student(Sno, Sdept, Dname)$ 中，$Sno \to Sdept$，$Sdept \to Dname$，则 $Dname$ 传递依赖于 $Sno$，不满足 3NF。

解决方法：分解为：
\begin{itemize}
  \item $Student(Sno, Sdept)$；
  \item $Dept(Sdept, Dname)$。
\end{itemize}

\subsection{BCNF}

\textcolor{red}{BCNF}（Boyce-Codd 范式）：在满足 3NF 的基础上，对于每一个非平凡函数依赖 $X \to Y$，$X$ 必须是候选键（超键）。

简单说：\textbf{凡是函数依赖的左侧，都必须是候选键}。

BCNF 是比 3NF 更严格的范式，3NF 的关系不一定满足 BCNF。

\subsection{范式之间的关系}

\begin{center}
\textbf{BCNF $\subset$ 3NF $\subset$ 2NF $\subset$ 1NF}
\end{center}

常见考试方式：给出关系模式和函数依赖集，判断属于哪个范式，并说明原因。

判断方法建议按照：\textbf{先判断主键，再看非主属性与主键的依赖关系，最后分析是否有传递依赖或左侧非候选键的情况}。

\subsection{函数依赖的公理系统（Armstrong 公理）}

考试有时会涉及 Armstrong 公理，共有三条基本规则：

\begin{enumerate}
  \item \textbf{自反律}：若 $Y \subseteq X$，则 $X \to Y$。
  \item \textbf{增广律}：若 $X \to Y$，则 $XZ \to YZ$。
  \item \textbf{传递律}：若 $X \to Y$，$Y \to Z$，则 $X \to Z$。
\end{enumerate}

通过这三条规则还可以推导出常用的合并律、分解律等。

\subsection{关系数据理论练习}

\begin{enumerate}
  \item 函数依赖 $X \to Y$ 的通俗含义是什么？
  \item 完全函数依赖和部分函数依赖的区别是什么？请举例说明。
  \item 设关系模式 $R(A, B, C)$，函数依赖集 $F = \{AB \to C, B \to A\}$，$R$ 满足第几范式？
  \item 2NF 和 3NF 分别消除了哪类函数依赖问题？
  \item 请判断：满足 3NF 的关系模式一定满足 BCNF。（\hspace{1cm}）
  \item 请判断：在关系模式中，若主键只有一个属性，则该模式一定满足 2NF。（\hspace{1cm}）
\end{enumerate}

\section{数据库设计}

\subsection{数据库设计的六个阶段}

数据库设计通常分为以下六个阶段：

\begin{enumerate}
  \item \textbf{需求分析}：调查用户需求，收集和分析数据，确定系统边界和目标。
  \item \textbf{概念结构设计}：将需求抽象成概念模型，通常使用 E-R 模型表达。
  \item \textbf{逻辑结构设计}：将概念模型转换为特定 DBMS 支持的数据模型，如关系模型。
  \item \textbf{物理结构设计}：确定数据在存储介质上的组织方式，包括索引、存储路径等。
  \item \textbf{数据库实施}：根据设计建立数据库，录入数据，测试应用程序。
  \item \textbf{数据库运行与维护}：长期运行中监控性能，进行调整和维护。
\end{enumerate}

考试中最常考的是前三个阶段，尤其是\textbf{概念结构设计}和\textbf{逻辑结构设计}。

\subsection{E-R 模型基础}

概念结构设计的核心工具是 \textbf{E-R 模型（实体-联系模型）}。

E-R 图中的三种基本元素：

\begin{itemize}
  \item \textbf{实体}：用矩形框表示，例如"学生"、"课程"；
  \item \textbf{属性}：用椭圆形表示，连接到对应实体；
  \item \textbf{联系}：用菱形框表示，连接相关实体，并标注联系类型（1:1、1:n、m:n）。
\end{itemize}

\subsubsection{联系的三种类型}

\begin{itemize}
  \item \textbf{一对一（1:1）}：一个实体实例只与另一方的一个实例相关联。例如：一个班级只有一个班长，一个班长只管理一个班级。
  \item \textbf{一对多（1:n）}：一个实体实例与另一方多个实例相关联。例如：一个老师讲授多门课程。
  \item \textbf{多对多（m:n）}：两方实体实例之间互为多关系。例如：一个学生选修多门课程，一门课程有多个学生选修。
\end{itemize}

\subsection{E-R 图转换为关系模式}

将 E-R 图转换为关系模式时，基本规则如下：

\begin{enumerate}
  \item 每个\textbf{实体}转换为一个关系模式，实体属性转为关系属性，实体的主标识符转为主键。
  \item \textbf{1:1 联系}：可将联系的属性合并到任一方关系中。
  \item \textbf{1:n 联系}：将联系及"一"端的主键加入"多"端关系中（即在"多"端加外键）。
  \item \textbf{m:n 联系}：联系需单独转换为一个关系模式，包含两端实体的主键（作为联合主键）和联系自身的属性。
\end{enumerate}

例如，学生与课程之间的"选修"（m:n 联系）需要单独建立选课关系：

\[
SC(Sno, Cno, Grade)
\]

其中 $(Sno, Cno)$ 为联合主键，$Sno$ 和 $Cno$ 分别是外键。

\subsection{数据库设计练习}

\begin{enumerate}
  \item 数据库设计分为哪六个主要阶段？
  \item 概念结构设计阶段的主要任务是什么？产出物是什么？
  \item E-R 图中矩形、椭圆和菱形分别表示什么？
  \item 三种联系类型（1:1、1:n、m:n）在向关系模式转换时分别采用什么策略？
  \item 若有实体"部门"（部门号、部门名）和"员工"（员工号、姓名），且一个部门有多个员工，应如何设计关系模式？
  \item 请判断：多对多联系在转换为关系模式时，不需要单独建立一个关系。（\hspace{1cm}）
\end{enumerate}

\section{数据库恢复与并发控制}

\subsection{事务}

\subsubsection{事务的定义}

\textcolor{red}{事务（Transaction）}是用户定义的一个数据库操作序列，这些操作要么全部执行，要么全部不执行，是一个不可分割的工作单位。

\subsubsection{事务的 ACID 属性}

事务具有四个基本属性，必须全部掌握：

\begin{enumerate}
  \item \textbf{原子性（Atomicity）}：事务中的所有操作要么全部成功，要么全部撤销，不存在部分执行的情况。
  \item \textbf{一致性（Consistency）}：事务执行前后，数据库必须从一个一致状态转移到另一个一致状态。
  \item \textbf{隔离性（Isolation）}：一个事务的执行不能被其他并发事务所干扰，各事务之间互相隔离。
  \item \textbf{持续性（Durability）}：事务一旦提交，其对数据库的改变就是永久性的，即使系统发生故障也不会丢失。
\end{enumerate}

记忆口诀：\textbf{原子一致隔离持久}（A-C-I-D）。

\subsection{故障类型}

数据库中可能遇到的故障主要有三类：

\begin{enumerate}
  \item \textbf{事务故障}：事务执行过程中由于逻辑错误或系统错误导致事务非正常终止。
  \item \textbf{系统故障}：整个数据库系统崩溃，内存中的数据丢失，但磁盘数据完好。
  \item \textbf{介质故障}（磁盘故障）：存储设备损坏，磁盘数据丢失，是最严重的一类故障。
\end{enumerate}

\subsection{恢复技术}

恢复技术的两个核心手段：

\begin{itemize}
  \item \textbf{日志（Log）}：记录每一次数据更新操作的内容（更新前后的值）；
  \item \textbf{备份（Backup）}：周期性地将数据库转储到安全介质。
\end{itemize}

恢复时的两种基本操作：

\begin{itemize}
  \item \textbf{Undo}（撤销）：将已执行但未提交的事务的操作撤销，还原数据；
  \item \textbf{Redo}（重做）：将已提交但尚未写入磁盘的事务操作重新执行。
\end{itemize}

\subsubsection{检查点技术}

检查点技术可以减少恢复时需要扫描的日志范围，提高恢复效率。发生故障后，只需从最近的检查点开始处理，而不是从头扫描整个日志。

\subsection{并发控制}

\subsubsection{并发操作引发的问题}

多个事务并发执行时可能出现以下三类问题：

\begin{enumerate}
  \item \textbf{丢失修改}：两个事务同时读取并修改同一数据，其中一个事务的修改被另一个覆盖而丢失。
  \item \textbf{不可重复读}：一个事务两次读取同一数据，其间另一个事务修改了该数据，导致两次读取结果不一致。
  \item \textbf{读脏数据}：一个事务读取了另一个事务未提交的修改数据，若后者回滚，则读到了无效数据。
\end{enumerate}

\subsubsection{封锁技术}

解决并发问题的主要机制是\textbf{封锁}：

\begin{itemize}
  \item \textbf{排他锁（X 锁、写锁）}：持锁事务可读可写，其他事务不能再加任何锁。
  \item \textbf{共享锁（S 锁、读锁）}：持锁事务只读，其他事务可以再加 S 锁，但不能加 X 锁。
\end{itemize}

\subsubsection{两段锁协议}

\textcolor{red}{两段锁协议（2PL）}是保证事务可串行化的常用协议：

\begin{itemize}
  \item \textbf{加锁阶段}：事务可以申请并获取锁，不能释放任何锁；
  \item \textbf{释放锁阶段}：事务可以释放锁，不能再申请新锁。
\end{itemize}

\subsection{恢复与并发控制练习}

\begin{enumerate}
  \item 事务的 ACID 属性分别是什么？请各用一句话说明。
  \item 数据库故障可分为哪三类？其中哪类最严重？
  \item Undo 和 Redo 分别在什么情况下使用？
  \item 并发事务可能引发哪三类问题？
  \item 排他锁和共享锁有什么区别？
  \item 请判断：只要事务提交，其对数据库的改变就一定是永久性的，这体现了事务的持续性。（\hspace{1cm}）
\end{enumerate}

\section{易混点总结}

\begin{enumerate}
  \item \textbf{安全性与完整性}：安全性防止\textbf{非法用户}访问，完整性防止\textbf{合法用户输入不合语义}的数据。
  \item \textbf{实体完整性与参照完整性}：前者约束\textbf{主键不为空}，后者约束\textbf{外键合法性}。
  \item \textbf{2NF 与 3NF}：2NF 消除部分函数依赖，3NF 消除传递函数依赖。
  \item \textbf{3NF 与 BCNF}：3NF 允许主属性之间存在函数依赖，BCNF 不允许。
  \item \textbf{概念结构设计与逻辑结构设计}：前者得到 E-R 图，后者得到关系模式。
  \item \textbf{m:n 联系转换}：必须单独建立一个关系，不能合并到任何一方。
  \item \textbf{Undo 与 Redo}：未提交事务用 Undo，已提交但未写磁盘的用 Redo。
  \item \textbf{X 锁与 S 锁}：X 锁排他，S 锁共享，多个事务可同时持有 S 锁但不能同时持有 X 锁。
\end{enumerate}

\section{本讲知识框架}

建议把本讲整理成下面这条主线：

\begin{center}
\begin{tabular}{|>{\centering\arraybackslash}m{2.2cm}|>{\centering\arraybackslash}m{2.8cm}|>{\centering\arraybackslash}m{3.2cm}|>{\centering\arraybackslash}m{3.2cm}|}
\hline
主题 & 核心问题 & 典型内容 & 高频考点 \\
\hline
安全性 & 谁能访问数据 & DAC、GRANT、REVOKE & 权限语句、DAC 与 MAC \\
\hline
完整性 & 数据合不合规 & 实体、参照、用户定义 & 三类完整性判断 \\
\hline
关系数据理论 & 模式设计合不合理 & 函数依赖、1NF～BCNF & 范式判断、模式分解 \\
\hline
数据库设计 & 如何系统地建库 & 六阶段、E-R 图、关系转换 & E-R 图绘制与转换 \\
\hline
恢复与并发 & 故障和并发如何处理 & ACID、日志、封锁 & ACID 属性、并发问题 \\
\hline
\end{tabular}
\end{center}

\section{典型例题}

\subsection{例题1：安全性与完整性的区别}

\textbf{题目：}数据库安全性和数据库完整性的主要区别是什么？

\textbf{分析：}两者都是保护数据库，但针对的对象不同。

\textbf{答案：}安全性是防止非法用户对数据库的非授权访问；完整性是防止合法用户向数据库中录入不合语义的数据。

\subsection{例题2：判断范式}

\textbf{题目：}设关系模式 $R(Sno, Cno, Sname, Grade)$，函数依赖集 $F = \{(Sno, Cno) \to Grade, Sno \to Sname\}$。$R$ 属于哪个范式？

\textbf{分析：}主键为 $(Sno, Cno)$，$Sname$ 仅依赖于 $Sno$，存在非主属性对主键的部分函数依赖。

\textbf{答案：}$R$ 满足 1NF，但不满足 2NF。

\subsection{例题3：模式分解}

\textbf{题目：}将上题中的 $R$ 分解为满足 2NF 的关系模式。

\textbf{分析：}将部分依赖的属性 $Sname$ 与其完全依赖的属性 $Sno$ 单独成表。

\textbf{答案：}分解为 $Student(Sno, Sname)$ 和 $SC(Sno, Cno, Grade)$。

\subsection{例题4：E-R 图转关系模式}

\textbf{题目：}若存在"学生"和"课程"两个实体，它们之间是多对多的"选修"联系，且选修联系含属性"成绩"，请将其转换为关系模式。

\textbf{分析：}m:n 联系需单独建立关系，包含两端主键及联系属性。

\textbf{答案：}
\begin{itemize}
  \item $Student(Sno, Sname, \ldots)$
  \item $Course(Cno, Cname, \ldots)$
  \item $SC(Sno, Cno, Grade)$，其中 $(Sno, Cno)$ 为主键。
\end{itemize}

\subsection{例题5：ACID 属性判断}

\textbf{题目：}银行转账操作中，从账户 A 扣款并向账户 B 入款，这两个操作必须同时成功或同时撤销，体现了事务的哪个属性？

\textbf{分析：}要么全部执行，要么全部不执行，是原子性的典型描述。

\textbf{答案：}原子性（Atomicity）。

\subsection{例题6：并发问题判断}

\textbf{题目：}事务 $T_1$ 读取数据 $x=100$，随即事务 $T_2$ 也读取 $x=100$，$T_1$ 将 $x$ 改为 200 并提交，$T_2$ 也将 $x$ 改为 120 并提交，最终 $x=120$ 而非 200。这属于哪类并发问题？

\textbf{分析：}$T_1$ 的修改被 $T_2$ 覆盖而丢失，属于丢失修改。

\textbf{答案：}丢失修改。

\section{课堂练习}

\subsection{基础练习}

\begin{enumerate}
  \item 数据库安全性控制的主要目标是什么？
  \item 三类数据库完整性约束分别是什么？各约束什么？
  \item 什么是函数依赖？请用一个生活例子说明。
  \item 数据库设计有哪六个阶段？各阶段的主要任务是什么？
  \item 事务的四个基本属性分别是什么？
\end{enumerate}

\subsection{判断与选择}

\begin{enumerate}
  \item 数据库安全性和完整性都是为了防止非法用户访问数据库。（\hspace{1cm}）
  \item 主键属性允许取 NULL 值。（\hspace{1cm}）
  \item 满足 3NF 的关系模式一定满足 BCNF。（\hspace{1cm}）
  \item E-R 图中的多对多联系转换为关系模式时，需要单独建立一个关系。（\hspace{1cm}）
  \item 事务提交后，即使系统故障，其对数据库的修改也应该被保留。（\hspace{1cm}）
  \item 两段锁协议可以解决所有并发控制问题。（\hspace{1cm}）
\end{enumerate}

\subsection{计算与分析}

\begin{enumerate}
  \item 设关系模式 $R(A, B, C, D)$，函数依赖集 $F = \{A \to B, B \to C, A \to D\}$，请指出 $R$ 的主键，并判断 $R$ 属于哪个范式。
  \item 将上题的 $R$ 分解为满足 3NF 的关系模式，并说明分解理由。
  \item 描述 Armstrong 公理的三条基本规则，并说明它们的用途。
  \item 说明排他锁和共享锁的兼容性规则：当一个事务持有 S 锁时，另一个事务能否再申请 X 锁？
  \item 请简述数据库备份与日志在故障恢复中各自的作用。
  \item Undo 和 Redo 操作分别针对哪类事务状态？
\end{enumerate}

\subsection{综合小题}

\textbf{题目：}某图书馆系统需要管理"读者"和"图书"两类实体，相关信息如下：

\begin{itemize}
  \item 读者：读者号、姓名、联系电话；
  \item 图书：书号、书名、作者；
  \item 一位读者可以借阅多本图书，一本图书可以被多位读者借阅；借阅关系含属性"借阅日期"。
\end{itemize}

请完成下面三个问题：

\begin{enumerate}
  \item 该系统中，读者与图书之间是什么类型的联系（1:1、1:n 还是 m:n）？
  \item 请根据描述设计对应的关系模式（给出每个关系及其属性，标注主键）。
  \item 若"借阅"关系为 $Borrow(Rno, Bno, BorrowDate)$，其中 $(Rno, Bno)$ 为主键，请说明 $BorrowDate$ 是否完全依赖于主键，并判断该关系满足几范式。
\end{enumerate}

\section{课后小测参考答案}

\begin{enumerate}
  \item 确保只有合法的用户以合法的方式访问数据。
  \item 实体完整性约束主键非空，参照完整性约束外键合法，用户定义完整性约束具体业务规则。
  \item 函数依赖 $X \to Y$ 是指当 $X$ 的值确定后，$Y$ 的值也唯一确定。例如，学号确定后姓名唯一确定。
  \item 需求分析、概念结构设计、逻辑结构设计、物理结构设计、数据库实施、数据库运行与维护。各阶段依次完成从需求收集、建模、转换到实现和维护的工作。
  \item 原子性、一致性、隔离性、持续性。
  \item 第1题错（完整性针对合法用户的不合规输入）；第2题错；第3题错；第4题对；第5题对；第6题错（两段锁不能完全避免死锁等问题）。
  \item $R$ 的主键为 $A$，存在传递依赖 $A \to B \to C$，故 $R$ 满足 2NF 但不满足 3NF。
  \item 分解为 $R_1(A, B, D)$ 和 $R_2(B, C)$，消除了传递依赖 $A \to B \to C$。
  \item Undo 针对已开始但未提交的事务，Redo 针对已提交但未写入磁盘的事务。
\end{enumerate}

综合小题参考：第 1 问为 m:n；第 2 问关系模式为 $Reader(\underline{Rno}, Rname, Tel)$、$Book(\underline{Bno}, Bname, Author)$、$Borrow(\underline{Rno, Bno}, BorrowDate)$；第 3 问 $BorrowDate$ 完全依赖于联合主键 $(Rno, Bno)$，该关系满足 3NF（非主属性只依赖于主键，无传递依赖）。

\subsection{分节练习参考答案}

\textbf{安全性练习参考：}

\begin{enumerate}
  \item DAC 由数据拥有者决定访问权限；MAC 由系统按安全标签强制执行，用户不能干预。
  \item 授予 user1 对 Student 表的查询（SELECT）和插入（INSERT）权限。
  \item 允许被授权用户将该权限再转授给其他用户。
  \item \texttt{REVOKE SELECT ON Student FROM user1;}
  \item 错。自主存取控制中，权限由数据拥有者或授权者决定。
  \item 错。\texttt{CASCADE} 会级联回收该用户转授出去的权限。
\end{enumerate}

\textbf{完整性练习参考：}

\begin{enumerate}
  \item 实体完整性要求主键不为空，因为主键负责唯一标识元组，若为空则无法识别。
  \item 参照完整性约束的是外键属性。
  \item 外键取值要么等于被参照关系中某个主键值，要么（若允许）为 NULL。
  \item 用户定义完整性是针对具体业务的约束，前两类是关系模型本身要求的通用约束。
  \item \texttt{CHECK} 用于定义属性取值范围，例如 \texttt{CHECK (Grade >= 0 AND Grade <= 100)}。
  \item 错。若业务允许，外键可以取 NULL 值。
\end{enumerate}

\textbf{关系数据理论练习参考：}

\begin{enumerate}
  \item $X$ 的值唯一确定后，$Y$ 的值也随之唯一确定。
  \item 完全函数依赖：$X$ 的任何真子集都不能单独确定 $Y$；部分函数依赖：存在 $X$ 的真子集可以确定 $Y$。例如 $(Sno, Cno) \to Grade$ 是完全依赖，$Sno \to Sname$ 是对 $(Sno, Cno)$ 的部分依赖。
  \item $AB$ 为主键；$B \to A$ 说明 $B$ 也可以确定 $A$，即 $A$ 部分依赖于主键 $AB$，$R$ 满足 1NF 但不满足 2NF。
  \item 2NF 消除部分函数依赖，3NF 消除传递函数依赖。
  \item 错。3NF 不一定满足 BCNF。
  \item 对。单属性主键不存在部分依赖，因此至少满足 2NF。
\end{enumerate}

\textbf{数据库设计练习参考：}

\begin{enumerate}
  \item 需求分析、概念结构设计、逻辑结构设计、物理结构设计、数据库实施、运行与维护。
  \item 主要任务是将需求抽象成概念模型；产出物是 E-R 图。
  \item 矩形表示实体，椭圆表示属性，菱形表示联系。
  \item 1:1 联系可合并入任一方；1:n 联系在多端加外键；m:n 联系需单独建立关系。
  \item $Dept(\underline{Dno}, Dname)$，$Employee(\underline{Eno}, Ename, Dno)$，其中 $Dno$ 为外键。
  \item 错。m:n 联系必须单独建立关系，以保存两端主键及联系属性。
\end{enumerate}

\textbf{恢复与并发控制练习参考：}

\begin{enumerate}
  \item 原子性：全做或全不做；一致性：执行前后数据库保持一致状态；隔离性：事务互不干扰；持续性：提交后永久生效。
  \item 事务故障、系统故障、介质故障；介质故障最严重。
  \item Undo 用于撤销未提交事务的操作；Redo 用于重做已提交但未写磁盘的操作。
  \item 丢失修改、不可重复读、读脏数据。
  \item X 锁排他，只有持锁者可读写；S 锁共享，多个事务可同时持有 S 锁，但有 S 锁时不能加 X 锁。
  \item 对。这体现了事务的持续性（Durability）。
\end{enumerate}

\section{本讲小结}

本讲是数据库部分的收尾，涵盖了从\textbf{数据保护（安全性、完整性）}到\textbf{模式优化（规范化理论）}再到\textbf{系统设计（数据库设计）}以及\textbf{系统可靠性（恢复与并发控制）}四个维度。

这四个维度共同支撑起一个完整的关系数据库系统：关系模型提供数据表示方式，SQL 提供操作手段，安全性和完整性约束保证数据合法可信，规范化理论保证模式设计合理，数据库设计流程保证从需求到实现的系统性，恢复与并发控制保证系统稳定运行。

至此，数据库原理部分的核心内容已全部覆盖。建议在考前重点回顾：\textbf{三类完整性、范式判断方法、E-R 图转换规则、事务 ACID 属性和三类并发问题}。

\end{document}
